%% toc: P15 162 162.2,162.6,162.10
\chapter{Calibrating Wald's Boundaries}
\label{sec:p15-w02c-p15-020-calibration}

\Cref{sec:p15-sequential-evidence} constructs Wald's probability
ratio test and proves its exact error inequalities. This chapter
adds the parts of the 1945 presentation that concern how the
boundaries are chosen and judged: the criterion against which the
rule is justified, the reason the posterior levels exceed one half,
what the usual boundaries \((1-\beta)/\alpha\) and \(\beta/(1-\alpha)\)
do and do not guarantee, when the nonzero exponential root used for
operating characteristics exists, and the source's dependent Gaussian
example \cite[Sections 2.2--3.4]{Wald1945SequentialTests}.
The notation is that of \cref{def:p15-se-model,def:p15-se-procedure,def:p15-se-sprt}:
\(L_m=p_{1m}/p_{0m}\), the test stops at the first \(m\) with
\(L_m\notin(B,A)\), and \(\alpha',\beta'\) are its actual errors.
For a sequential test \(S\), the source writes \(E_i(n\mid S)\) for
its expected number of observations under \(H_i\); here that number
is the stopping time \(T\), so \(E_i(n\mid S)=E_iT\).

\section{The criterion and the posterior levels}
\label{sec:p15-w02c-p15-020-criteria}

\begin{hbdefinition}{def:p15-w02c-p15-020-optimum}{Optimum tests and the two cost criteria}
Strength, \emph{better} and \emph{admissible} are as in
\cref{def:p15-se-efficiency}. Fix a sequential test \(S\) and let
\(\mathcal S\) be the class of sequential tests of strength equal to
that of \(S\). A test \(S\in\mathcal S\) is an \emph{optimum test}
if it is admissible and
\[
 \max\{E_0(n\mid S),E_1(n\mid S)\}\leq
 \max\{E_0(n\mid S'),E_1(n\mid S')\}
 \quad\text{for every }S'\in\mathcal S .
\]
For fixed weights \(w_0,w_1>0\), the \emph{weighted criterion}
minimizes \(w_0E_0(n\mid S)+w_1E_1(n\mid S)\) over \(\mathcal S\).

The source's reason for the choice is as follows. If it were known how
often \(H_0\) and \(H_1\) hold in the long run, the natural target
would be the average of \(E_0(n)\) and \(E_1(n)\) weighted by those
frequencies. When no such knowledge is available, as the source
judges to be usual in practice, minimizing the larger of the two
expectations is preferred, and this is the definition the source
adopts \cite[Section 2.2, p.~124]{Wald1945SequentialTests}.
The source does not prove that a test minimizing both expectations at
once exists; it states that its probability ratio test very nearly
minimizes both, conjectures exact minimization, and treats the test as
optimum for practical purposes
\cite[Section 2.2, p.~124, footnote 4; Section 3.1, p.~127, footnote 8]{Wald1945SequentialTests}.
The proved part of that claim in this handbook is
\cref{thm:p15-se-no-overshoot}.
\end{hbdefinition}

\begin{hbtheorem}{thm:p15-w02c-p15-020-criteria-coincide}{proved=annex:p15-w02c-p15-020-criteria-coincide}{The criteria agree when a simultaneous minimizer exists}
Suppose some \(S^*\in\mathcal S\) has finite expectations and
satisfies \(E_i(n\mid S^*)\leq E_i(n\mid S')\) for \(i=0,1\) and every
\(S'\in\mathcal S\). For \(S\in\mathcal S\) and any fixed
\(w_0,w_1>0\), the following are equivalent:
\begin{enumerate}
\item \(S\) is admissible;
\item \(E_0(n\mid S)=E_0(n\mid S^*)\) and \(E_1(n\mid S)=E_1(n\mid S^*)\);
\item \(S\) is an optimum test;
\item \(S\) minimizes \(w_0E_0(n\mid S)+w_1E_1(n\mid S)\) over \(\mathcal S\).
\end{enumerate}
In particular \(S^*\) is an optimum test. This is the source's remark
that its candidate definitions coincide when both expectations can be
minimized together \cite[Section 2.2, p.~124]{Wald1945SequentialTests}.
Without admissibility the maximum criterion alone would not suffice:
a test with \(E_0(n\mid S)>E_0(n\mid S^*)\) and
\(E_1(n\mid S)=E_1(n\mid S^*)\geq E_0(n\mid S)\) has the same maximum
as \(S^*\).
\end{hbtheorem}

\begin{hbtheorem}{thm:p15-w02c-p15-020-posterior-majority}{proved=annex:p15-w02c-p15-020-posterior-majority}{Why both posterior levels exceed one half}
In the posterior rule of \cref{thm:p15-se-posterior}, accept \(H_1\)
when \(g_{1m}\geq d_1\) and \(H_0\) when \(g_{0m}\geq d_0\), where
\(g_{im}\) is the posterior probability of \(H_i\) after \(m\)
observations.
\begin{enumerate}
\item If \(d_0,d_1>1/2\), then whenever the rule accepts \(H_i\),
 \(g_{im}>g_{1-i,m}\): the accepted hypothesis has the larger
 posterior probability.
\item If \(d_1\leq1/2\), this can fail with positive probability.
 Let \(g_0=g_1=1/2\), let \(f_0\) and \(f_1\) be the Bernoulli laws
 with success probabilities \(3/4\) and \(1/2\), and let
 \(d_1=2/5\), \(d_0=3/4\). After a first observation equal to \(1\),
 which has probability \(3/4\) under \(H_0\), one has
 \(g_{11}=2/5\geq d_1\) and \(g_{01}=3/5<d_0\), so the rule accepts
 \(H_1\) although \(H_0\) has the larger posterior probability.
 The case \(d_0\leq1/2\) is the same with the hypotheses exchanged.
\end{enumerate}
This is the reason the source gives for imposing
\(d_0,d_1>1/2\) \cite[Section 3.1, p.~126, footnote 7]{Wald1945SequentialTests}.
\end{hbtheorem}

\begin{hbtheorem}{thm:p15-w02c-p15-020-fixed-size}{proved=annex:p15-w02c-p15-020-fixed-size}{When the number of observations is fixed in advance}
Let \(N\geq1\), and let a sequential procedure, possibly randomized as
in \cref{def:p15-se-procedure}, satisfy \(T\leq N\) almost surely
under both laws, with actual errors \(\alpha',\beta'\).
Then a most powerful test of size \(\alpha'\) based on all \(N\)
observations, as in \cref{thm:p15-neyman-pearson} applied to the
joint densities \(p_{0N},p_{1N}\), has second-kind error at most
\(\beta'\).
Hence, when exactly \(N\) observations will be available whatever
happens, no sequential use of them improves on the fixed most powerful
test in error probabilities. This is the source's statement that in
that situation the current most powerful test is preferable, since the
advantage of the sequential test lies only in a smaller expected
number of observations \cite[Introduction, p.~119]{Wald1945SequentialTests}.
\end{hbtheorem}

\section{The usual boundaries and what they guarantee}
\label{sec:p15-w02c-p15-020-nominal}

For prescribed \(\alpha,\beta\in(0,1)\) with \(\alpha+\beta<1\), the
source writes
\[
 a(\alpha,\beta)=\frac{1-\beta}{\alpha},\qquad
 b(\alpha,\beta)=\frac{\beta}{1-\alpha},
\]
and \(A(\alpha,\beta),B(\alpha,\beta)\) for boundaries whose actual
errors are exactly \(\alpha\) and \(\beta\)
\cite[Section 3.3, p.~131]{Wald1945SequentialTests}.
The condition \(\alpha+\beta<1\) is equivalent to
\(b(\alpha,\beta)<a(\alpha,\beta)\) and implies
\(b(\alpha,\beta)<1<a(\alpha,\beta)\)
\cite[Section 3.4, p.~134]{Wald1945SequentialTests}.

The test with boundaries \(a(\alpha,\beta),b(\alpha,\beta)\) needs no
distribution theory to be carried out. The boundaries depend on
\(\alpha\) and \(\beta\) alone, and the ratio \(L_m\) is evaluated from
the observed data and the two specified densities. The guarantees of
\cref{thm:p15-se-nominal} and of the next two theorems hold for every
pair of distinct laws, with no calculation of the distribution of a
test statistic. A distributional calculation is needed only for
statements about the number of observations the test will require,
such as \cref{thm:p15-se-mean} and \cref{sec:p15-se-time-transforms}.
The source contrasts this with the fixed-sample procedure, which needs
the distribution of its test statistic
\cite[Introduction, p.~119; Section 3.3, p.~133]{Wald1945SequentialTests}.
It adds that the distribution of the number of observations is of
secondary importance, so long as it is known that the sequential test
saves observations on average \cite[Section 3.3, p.~133]{Wald1945SequentialTests}.
That ranking is the source's practical judgement, not a theorem. What
can be proved is the average saving in a stated regime
(\cref{thm:p15-w02c-p15-020-small-shift}); the law of the number of
observations is approximated in \cref{sec:p15-se-time-transforms} and
\cref{thm:p15-w02c-p15-020-one-sided}, and a forced deadline changes
the errors (\cref{thm:p15-w02c-p15-020-truncation-exact}).

\begin{hbtheorem}{thm:p15-w02c-p15-020-one-boundary}{proved=annex:p15-w02c-p15-020-one-boundary}{Moving one boundary at a time}
Let \(0<B<A_1\leq A_2<\infty\) and compare the ratio tests with
boundaries \((B,A_1)\) and \((B,A_2)\). On every sample path the
second test does not stop before the first, it accepts \(H_1\) only on
paths where the first does, and it accepts \(H_0\) on every path where
the first does. Hence, under any law of the observations, raising the
upper boundary with the lower one fixed does not increase \(\alpha'\)
and does not decrease \(\beta'\).
Symmetrically, for \(0<B_2\leq B_1<A<\infty\), lowering the lower
boundary with the upper one fixed does not increase \(\beta'\) and does
not decrease \(\alpha'\); again the stopping time does not decrease.
The changes are weak: if no attainable value of the ratio lies in
\([A_1,A_2)\), respectively in \((B_2,B_1]\), the two tests coincide.
The source states the directions as a lower \(\alpha\) with a slightly
higher \(\beta\), and conversely
\cite[Section 3.3, p.~131]{Wald1945SequentialTests}.
\end{hbtheorem}

\begin{hbtheorem}{thm:p15-w02c-p15-020-nominal-sharp}{proved=annex:p15-w02c-p15-020-nominal-sharp}{At most one error can exceed its target, and when the sum inequality is strict}
Let \(P_0\neq P_1\), let \(\alpha,\beta\in(0,1)\) with
\(\alpha+\beta<1\), and use \(A=a(\alpha,\beta)\),
\(B=b(\alpha,\beta)\). Then:
\begin{enumerate}
\item at least one of \(\alpha'\leq\alpha\) and \(\beta'\leq\beta\)
 holds;
\item the following are equivalent:
 \(\alpha'+\beta'=\alpha+\beta\); \(\alpha'=\alpha\) and
 \(\beta'=\beta\); the test has no overshoot at either boundary, that
 is,
 \[
  P_1(Q_1\cap\{L_T\neq A\})=0,\qquad P_0(Q_0\cap\{L_T\neq B\})=0;
 \]
\item the possible excesses are bounded by
 \[
  \alpha'-\alpha\leq\frac{\alpha\beta}{1-\beta},\qquad
  \beta'-\beta\leq\frac{\alpha\beta}{1-\alpha}.
 \]
\end{enumerate}
Consequently, whenever the test with boundaries
\(a(\alpha,\beta),b(\alpha,\beta)\) does not itself have errors exactly
\(\alpha\) and \(\beta\), one has \(\alpha'+\beta'<\alpha+\beta\).
For \(\alpha=\beta=1/20\) each excess in (3) is at most \(1/380\).
\end{hbtheorem}

\begin{hbtheorem}{thm:p15-w02c-p15-020-exact-pair}{proved=annex:p15-w02c-p15-020-exact-pair}{The usual boundaries are wider than any exact pair}
Let \(P_0\neq P_1\) and \(\alpha,\beta\in(0,1)\) with
\(\alpha+\beta<1\). If boundaries \(0<B^*<A^*<\infty\) give a ratio
test with actual errors exactly \(\alpha\) and \(\beta\), then
\[
 A^*\leq a(\alpha,\beta),\qquad B^*\geq b(\alpha,\beta),
\]
the test with boundaries \(a(\alpha,\beta),b(\alpha,\beta)\) stops no
earlier than the exact one on every sample path, and so
\(E_0T\) and \(E_1T\) are not decreased by the substitution.
Exact boundaries need not be unique: two pairs that stop and decide
identically on every path have the same errors. So the source's
condition \(a(\alpha,\beta)>A(\alpha,\beta)\) or
\(b(\alpha,\beta)<B(\alpha,\beta)\) for a strict sum inequality is
made precise by the last statement of
\cref{thm:p15-w02c-p15-020-nominal-sharp}, which refers only to the
errors of the test actually used.
These are the source's (3.24)--(3.25) and its statements that the
substitution can only increase the expected numbers of observations and
gains strength when the inequalities are strict
\cite[Section 3.3, pp.~131 and 133; Section 4.3, p.~146]{Wald1945SequentialTests}.
\end{hbtheorem}

\begin{hbwarning}{def:p15-w02c-p15-020-why-inequalities}{The source's account of the gap between the exact and the usual boundaries}
The source explains why (3.24)--(3.25) are inequalities rather than
equalities: the test may stop with the ratio strictly beyond a
boundary. Were the ratio exactly at a boundary on stopping, the exact
and the usual boundaries would agree; this is item (2) of
\cref{thm:p15-w02c-p15-020-nominal-sharp}. The source then makes three
further statements that this handbook reports as its judgements rather
than proves in general.
First, when \(f_1\) is close to \(f_0\), one further observation
changes the ratio only slightly, so the final ratio lies only slightly
beyond the boundary and the exact boundaries are nearly the usual ones.
Second, if the number of observations could vary continuously, the
ratio would reach the boundary exactly and the exact boundaries would
equal the usual ones; the inequalities are due only to the integer
number of observations.
Third, for small targets, say below \(0.05\), any excess of
\(\alpha'\) over \(\alpha\) or of \(\beta'\) over \(\beta\) is small
enough to neglect, and the source points to its Section 3.4 bounds as
showing that the excess is extremely small
\cite[Section 3.3, pp.~132--133]{Wald1945SequentialTests}.
Item (3) of \cref{thm:p15-w02c-p15-020-nominal-sharp} is a proved
bound of this kind; the Section 3.4 bounds appear here as
\cref{thm:p15-se-oc}. The source's remark that the only possible
disadvantage of the usual boundaries is a slight increase in the
expected number of observations is likewise a judgement: its direction
is \cref{thm:p15-w02c-p15-020-exact-pair}, while its size is not
bounded in general here. \Cref{thm:p15-se-nominal} shows that an
individual error can exceed its target when the laws are far apart.
\end{hbwarning}

\section{The nonzero exponential root}
\label{sec:p15-w02c-p15-020-root}

The operating-characteristic calculation of
\cref{def:p15-se-exponential-root,thm:p15-se-oc} assumes a real
\(h\neq0\) with \(\varphi(h)=Ee^{hz}=1\) for the score law \(z\) under
the true hypothesis. The source states that, under very mild
restrictions on the distribution of \(z\), exactly one such \(h\)
exists; it attributes this to its reference [4] and does not list the
restrictions in this paper. It also assumes that the variance of
\(z\) is not zero, and states that \(h=1\) when \(H_0\) is true and
\(h=-1\) when \(H_1\) is true
\cite[Section 3.4, pp.~134--135]{Wald1945SequentialTests}.
The next three results give a sufficient set of restrictions with a
full proof, prove the two values under the specified laws without any
moment assumption, and show that some restriction is necessary.

\begin{hbtheorem}{thm:p15-w02c-p15-020-root-exists}{proved=annex:p15-w02c-p15-020-root-exists}{Existence and uniqueness of the nonzero root}
Let \(z\) be a real random variable with
\begin{enumerate}
\item \(\varphi(t)=Ee^{tz}<\infty\) for every real \(t\);
\item \(P(z>0)>0\) and \(P(z<0)>0\).
\end{enumerate}
Then \(E|z|<\infty\). If \(Ez<0\), there is exactly one real
\(h\neq0\) with \(\varphi(h)=1\), and \(h>0\); if \(Ez>0\), there is
exactly one, and \(h<0\). If \(Ez=0\), then \(\varphi(t)>1\) for every
\(t\neq0\), so no nonzero root exists.
\end{hbtheorem}

\begin{hbtheorem}{thm:p15-w02c-p15-020-root-simple-laws}{proved=annex:p15-w02c-p15-020-root-simple-laws}{The root is one under the null law and minus one under the alternative}
Let \(P_0,P_1\) be equivalent and distinct, and let
\(z=\log(f_1(X)/f_0(X))\). Under \(H_0\), the function
\(t\mapsto E_0e^{tz}\) is finite on \([0,1]\) and \(t=1\) is the only
real \(t\neq0\) at which it equals one. Under \(H_1\), \(t=-1\) is the
only real \(t\neq0\) at which \(E_1e^{tz}\) equals one.
No moment of \(z\) is assumed.
\end{hbtheorem}

\begin{hbtheorem}{thm:p15-w02c-p15-020-root-failure}{proved=annex:p15-w02c-p15-020-root-failure}{A nonconstant score law with no nonzero root}
Let \(W\geq0\) have \(P(W>w)=(1+w)^{-2}\) for \(w\geq0\), and put
\(z=W-2\). Then \(Ez=-1\), \(P(z>0)=1/9\), \(P(z<0)=8/9\), and
\(\varphi(t)=\infty\) for every \(t>0\), while \(\varphi(t)>1\) for
every \(t<0\). Thus \(\varphi(h)=1\) only at \(h=0\). The finiteness
restriction (1) of \cref{thm:p15-w02c-p15-020-root-exists} cannot
simply be dropped, and \cref{thm:p15-se-oc} has no root to use for
such a score law.
\end{hbtheorem}

\section{A dependent Gaussian example}
\label{sec:p15-w02c-p15-020-autoregressive}

The source observes that independence was used only to prove that the
test terminates with probability one, its equation (3.11). It states
that termination can be shown for dependent observations under very
general conditions, and that (3.18)--(3.19) then hold for the ratio of
joint densities \cite[Section 3.2, pp.~128 and 130]{Wald1945SequentialTests}.
The general statement is \cref{thm:p15-se-dependent}, which assumes
termination; the source does not specify the general conditions, and
they are not formulated here. The source's example follows, with its
termination proved.

\begin{hbtheorem}{thm:p15-w02c-p15-020-autoregressive}{proved=annex:p15-w02c-p15-020-autoregressive}{An autoregressive pair of hypotheses}
Let \(\lambda_0\neq\lambda_1\) lie in \((0,1)\). Under \(H_i\), the
observations \(x_1,x_2,\ldots\) have joint densities
\[
 p_{im}(x_1,\ldots,x_m)=(2\pi)^{-m/2}
 \exp\Bigl\{-\tfrac12x_1^2-\tfrac12\sum_{j=2}^m(x_j-\lambda_ix_{j-1})^2\Bigr\},
\]
so that \(x_1\) and the residuals \(x_j-\lambda_ix_{j-1}\), \(j\geq2\),
are independent standard normal variables. Then:
\begin{enumerate}
\item the densities are consistent: integrating \(p_{i,m+1}\) over its
 last coordinate gives \(p_{im}\);
\item with \(\delta=\lambda_1-\lambda_0\) and
 \(\bar\lambda=(\lambda_0+\lambda_1)/2\),
 \[
  \log\frac{p_{1m}}{p_{0m}}
  =\delta\sum_{j=2}^mx_{j-1}(x_j-\bar\lambda x_{j-1}),
 \]
 which tends to \(+\infty\) almost surely under \(H_1\) and to
 \(-\infty\) almost surely under \(H_0\);
\item for every \(0<B<A<\infty\), the ratio test applied to
 \(p_{1m}/p_{0m}\) terminates almost surely under both hypotheses, and
 its actual errors satisfy
 \[
  A\alpha'\leq1-\beta',\qquad \beta'\leq B(1-\alpha').
 \]
\end{enumerate}
These are the source's (3.18)--(3.19) for this example, written
without division
\cite[Section 3.2, p.~130]{Wald1945SequentialTests}. Replacing the
joint ratio by a product of one-dimensional marginal ratios would not
give this test.
\end{hbtheorem}
\section{When the boundary approximations are accurate}
\label{sec:p15-w02c-p15-020-regimes}

In this section \(a=\log A\) and \(b=\log B\) are the logarithmic
boundaries of \cref{def:p15-se-sprt}, with \(0<B<1<A\), and the scores
\(z_1,z_2,\ldots\) are independent copies of a finite score \(z\) under
the law in force. The source replaces the conditional terminal
quantities \(E^*=E(e^{hZ_T}\mid Q_0)\) and \(E^{**}=E(e^{hZ_T}\mid Q_1)\)
by \(B^h\) and \(A^h\), and the conditional terminal sums by \(b\) and
\(a\), when \(|Ez|\) and the variance or standard deviation of \(z\)
are small. It states that \(\bar\gamma\) is then a good approximation
to \(\gamma\) and that \(\bar\gamma-\gamma\) tends to zero when the
mean and the variance of \(z\) tend to zero
\cite[Section 3.4, p.~135; Section 4.1, p.~143]{Wald1945SequentialTests}.
The next two results separate what is true in that statement: a uniform
bound on the size of one score suffices, while small mean and variance
alone do not.

\begin{hbtheorem}{thm:p15-w02c-p15-020-gamma-bar}{proved=annex:p15-w02c-p15-020-gamma-bar}{Bounded scores make the approximation to the operating characteristic accurate}
Let \(|z|\leq c\) almost surely for a constant \(c>0\), let \(h\neq0\)
satisfy \(Ee^{hz}=1\), and let \(T<\infty\) almost surely. Put
\(s=e^{|h|c}\). If \(h>0\),
\[
 \frac{1-B^h}{sA^h-B^h}\leq\gamma\leq\frac{1-B^h/s}{A^h-B^h/s};
\]
if \(h<0\), the same holds for \(1-\gamma\) with \(A^h\) and \(B^h\)
exchanged. As \(s\to1\) the bounds tend to \(\bar\gamma=(1-B^h)/(A^h-B^h)\),
respectively to \(1-\bar\gamma\), uniformly for \(h\) in any compact
subset of \(\mathbb R\setminus\{0\}\). Consequently, for score laws
\(z_k\) with roots \(h_k\) in such a compact set and bounds
\(|z_k|\leq c_k\to0\), \(\gamma_k-\bar\gamma_k\to0\); the mean and the
variance of \(z_k\) then tend to zero as well. For the two specified
hypotheses, where \(h=1\) under \(H_0\) and \(h=-1\) under \(H_1\)
(\cref{thm:p15-w02c-p15-020-root-simple-laws}), the condition is that
\(\sup|\log(f_1/f_0)|\) be small, one exact form of the source's
phrase that \(f_1\) is near \(f_0\).
\end{hbtheorem}

\begin{hbtheorem}{thm:p15-w02c-p15-020-gamma-bar-failure}{proved=annex:p15-w02c-p15-020-gamma-bar-failure}{Small mean and variance do not suffice}
Let \(a=1/4\) and \(b=-1/4\), so \(A=e^{1/4}\) and \(B=e^{-1/4}\). For
\(0<p<1/e\), let \(z\) take the value \(1\) with probability \(p\) and
the value \(-\epsilon_p\) with probability \(1-p\), where
\(\epsilon_p=\log\{(1-p)/(1-pe)\}>0\). Then \(Ee^{z}=1\), so \(h=1\) is
the unique nonzero root; \(Ez<0\); and \(Ez\to0\) and
\(\operatorname{Var}z\to0\) as \(p\to0\). The test terminates almost
surely, every upper exit overshoots \(a\) by more than \(1/2\), and
\[
 \gamma\leq\frac{1-Be^{-\epsilon_p}}{Ae^{1/2}-Be^{-\epsilon_p}}
 \longrightarrow\frac{1-B}{Ae^{1/2}-B}<\frac{1-B}{A-B}=\bar\gamma
 \qquad(p\to0).
\]
So \(\bar\gamma-\gamma\) stays bounded away from zero although the mean
and the variance of the score tend to zero.
\end{hbtheorem}

\begin{hbtheorem}{thm:p15-w02c-p15-020-lattice}{proved=annex:p15-w02c-p15-020-lattice}{Lattice approximations of small span}
For a span \(d>0\) put \(z^-=d\lfloor z/d\rfloor\) and
\(z^+=d\lceil z/d\rceil\), computed from the same observations, and let
\(Z_m^\pm\), \(T^\pm\) and \(\gamma^\pm_d\) be the sums, the exit times
from \((b,a)\) under the rule of \cref{def:p15-se-sprt}, and the
probabilities of an upper exit of the two lattice walks.
\begin{enumerate}
\item For every \(d\), \(\gamma^-_d\leq\gamma\leq\gamma^+_d\), where
 \(\gamma\) is the probability of an upper exit of the original walk.
 When \(z\) is bounded, \(z^\pm/d\) takes finitely many integer values,
 so the lattice quantities are absorption quantities of a finite-state
 walk, computable exactly as in \cref{thm:p15-se-finite-recursion}
 whenever \(z^\pm/d\) takes both signs.
\item If \(T<\infty\) almost surely and \(P(Z_m\in\{a,b\})=0\) for every
 \(m\geq1\), which holds whenever the law of \(z\) has no atoms, then
 almost surely \(T^\pm=T\), with the same decision, for all small enough
 \(d\); hence \(\gamma^\pm_d\to\gamma\) as \(d\to0\).
\item If in addition \(|z|\leq c\), then
 \(E(Z^\pm_{T^\pm};T^\pm<\infty)\to EZ_T\), and when \(Ez\neq0\) and
 \(ET<\infty\), \(ET=EZ_T/Ez\) is approximated by these lattice values
 divided by \(Ez\).
\end{enumerate}
These are exact forms of the source's statements that a score that is
not confined to the multiples of a constant can be approximated by one
that is, which approximates \(\gamma\) and \(EZ_n\) arbitrarily finely
\cite[Section 3.4, p.~136; Section 4.1, p.~143]{Wald1945SequentialTests}.
When the walk can land exactly on a boundary with positive probability,
item (1) still brackets \(\gamma\), but the bracket need not shrink.
For unbounded scores, convergence of the terminal means in item (3) is
not proved here.
\end{hbtheorem}

\begin{hbtheorem}{thm:p15-w02c-p15-020-small-shift}{proved=annex:p15-w02c-p15-020-small-shift}{The normal small-shift regime}
Fix \(\alpha,\beta\in(0,1)\) with \(\alpha+\beta<1\), use the boundaries
\(A=a(\alpha,\beta)\) and \(B=b(\alpha,\beta)\), so that
\(a=\log\{(1-\beta)/\alpha\}\) and \(b=\log\{\beta/(1-\alpha)\}\), and
for \(d>0\) let the hypotheses be \(N(\theta_0,1)\) and \(N(\theta_1,1)\)
with \(\theta_1-\theta_0=2d\). Write \(\alpha'_d,\beta'_d\) for the
actual errors, \(E_{i,d}T\) for the expected numbers of observations,
\(q_u=\Phi^{-1}(u)\), and \(n_d=(q_{1-\alpha}+q_{1-\beta})^2/(4d^2)\),
the source's \(n(\alpha,\beta)\) of (4.30) before rounding.
\begin{enumerate}
\item Under either hypothesis the overshoot constants of
 \cref{def:p15-se-overshoot} satisfy
 \[
  \max\{\xi,-\xi'\}\leq M_d=2d\Bigl\{d+\frac{\phi(d)}{1-\Phi(d)}\Bigr\},
 \]
 and
 \[
  \bigl|\,2d^2E_{1,d}T-\{(1-\beta'_d)a+\beta'_db\}\bigr|\leq M_d,
 \]
 \[
  \bigl|\,2d^2E_{0,d}T+\{(1-\alpha'_d)b+\alpha'_da\}\bigr|\leq M_d .
 \]
 For \(0<d\leq1\), \(M_d\leq2d(1+\phi(0)/(1-\Phi(1)))<8d\).
\item As \(d\to0\), \(\alpha'_d\to\alpha\) and \(\beta'_d\to\beta\).
\item As \(d\to0\),
 \[
  \frac{E_{1,d}T}{n_d}\to R_1(\alpha,\beta)=
  \frac{2\{(1-\beta)a+\beta b\}}{(q_{1-\alpha}+q_{1-\beta})^2},
 \]
 \[
  \frac{E_{0,d}T}{n_d}\to R_0(\alpha,\beta)=
  \frac{-2\{(1-\alpha)b+\alpha a\}}{(q_{1-\alpha}+q_{1-\beta})^2},
 \]
 and the same limits hold with \(n_d\) replaced by
 \(N_d=\lceil n_d\rceil\).
\item Let \(I_{i,d}\) be the numerator of the lower bound of
 \cref{thm:p15-se-information} under \(H_i\), evaluated at
 \(\alpha'_d,\beta'_d\), and \(\mu_i=E_iz\). Then
 \(1\leq|\mu_i|E_{i,d}T/I_{i,d}\to1\) for \(i=0,1\).
\end{enumerate}
Item (1) is the regime in which the source neglects the excess: the
error is of order \(d\), that is of the order of
\(|\theta_1-\theta_0|\) and of the standard deviation \(2d\) of \(z\),
against main terms that do not shrink
\cite[Section 4.1, p.~143; Section 4.3, p.~146]{Wald1945SequentialTests}.
Items (2) and (3) show that the overshoot-free values behind Table 1 are
the limits of the exact ones. Item (4) is the exact content of the
source's statement in Section 4.7 that small \(|Ez|\) and \(\sigma_z\)
make the expected excess small and the test very nearly optimum
\cite[Section 4.7, pp.~154 and 157]{Wald1945SequentialTests}: by
\cref{thm:p15-se-information}, every terminating test with the same
actual errors needs at least \(I_{i,d}/|\mu_i|\) observations on
average under \(H_i\), so the ratio of the expected number for this
test to that of any such competitor is at most a quantity tending to
one.
\end{hbtheorem}

\begin{hbwarning}{def:p15-w02c-p15-020-zero-excess}{When the excess is zero}
The source states that the excess over the boundaries is exactly zero
when \(z\) takes only the values \(d\) and \(-d\) and \(\log A\),
\(\log B\) are integral multiples of \(d\), and that it is not
identically zero in any other case
\cite[Section 4.7, p.~154]{Wald1945SequentialTests}. The first part is
the nearest-neighbour case of \cref{thm:p15-se-no-overshoot}. The second
part is too strong for a fixed pair of boundaries. For \(d>0\), let
\(X\) be Bernoulli with success probability
\(p_0=(e^{2d}-1)/(e^{3d}-1)\) under \(H_0\) and \(p_1=e^dp_0\) under
\(H_1\). Then \(z=d\) when \(X=1\) and \(z=-2d\) when \(X=0\), and with
\(A=e^d\), \(B=e^{-2d}\) the test stops at the first observation exactly
on a boundary. It has no excess, so by \cref{thm:p15-se-no-overshoot} it
is exactly optimum, although \(z\) is not of the form \(\pm d\). The
verification is the last paragraph of the proof of
\cref{thm:p15-w02c-p15-020-small-shift}.
\end{hbwarning}

\section{Moments and the law of the number of observations}
\label{sec:p15-w02c-p15-020-moments}

\begin{hbtheorem}{thm:p15-w02c-p15-020-moment-system}{proved=annex:p15-w02c-p15-020-moment-system}{The linear equations for the moments under the boundary approximation}
Let \(z\) be nonconstant, let \(\varphi(t)=Ee^{tz}\) be finite on an
open interval containing \(0\) and a root \(h\neq0\) of
\(\varphi(h)=1\), and let \(\mu=Ez\neq0\). Put
\(\ell(t)=-\log\varphi(t)\), \(\mu_h=\varphi'(h)\) and
\(\bar\gamma=(1-B^h)/(A^h-B^h)\in(0,1)\) as in \cref{thm:p15-se-oc}.
Suppose functions \(\widehat\Psi_+,\widehat\Psi_-\), \(r\) times
differentiable at zero, satisfy the source's approximate identity
(4.47),
\[
 \bar\gamma A^t\widehat\Psi_+(\ell(t))+(1-\bar\gamma)B^t\widehat\Psi_-(\ell(t))=1,
\]
for \(t\) near \(0\) and near \(h\). Then
\(\widehat\Psi_+(0)=\widehat\Psi_-(0)=1\), and differentiating \(j\)
times, \(1\leq j\leq r\), at \(t=0\) and at \(t=h\) gives \(2r\) linear
equations in the \(2r\) numbers \(\widehat\Psi_\pm^{(j)}(0)\). They are
uniquely solvable order by order: the equations of order \(j\) form a
\(2\times2\) system for \(\widehat\Psi_+^{(j)}(0),\widehat\Psi_-^{(j)}(0)\)
with determinant
\(\bar\gamma(1-\bar\gamma)\mu^j\mu_h^j(B^h-A^h)\neq0\), whose other
terms involve only lower orders. For \(j=1\), with
\(u=\bar\gamma\widehat\Psi_+'(0)\) and \(v=(1-\bar\gamma)\widehat\Psi_-'(0)\),
the equations are the source's (4.49) and (4.50):
\[
 \mu(u+v)=\bar\gamma a+(1-\bar\gamma)b,\qquad
 \mu_h(A^hu+B^hv)=\bar\gamma aA^h+(1-\bar\gamma)bB^h,
\]
so the approximate mean \(u+v\) is the overshoot-free value of
\cref{thm:p15-se-mean}, and
\[
 u=\frac{\{\bar\gamma a+(1-\bar\gamma)b\}B^h/\mu
   -\{\bar\gamma aA^h+(1-\bar\gamma)bB^h\}/\mu_h}{B^h-A^h},
\]
\[
 v=\frac{\bar\gamma a+(1-\bar\gamma)b}{\mu}-u .
\]
This is the source's derivation of the moments of \(n\) from (4.35)
through (4.47)--(4.50)
\cite[Section 4.4, pp.~149--150]{Wald1945SequentialTests}. The numbers
obtained are moments of the approximating transforms of
\cref{def:p15-se-transform-approximation}; they are moments of \(T\)
only to the extent that the boundary approximation holds.
\end{hbtheorem}

\begin{hbtheorem}{thm:p15-w02c-p15-020-one-sided}{proved=annex:p15-w02c-p15-020-one-sided}{The one-sided cases and the effect of a rarely reached lower boundary}
\begin{enumerate}
\item (The case \(B>0\), \(A=\infty\), \(Ez<0\).) Stopping at the first
 \(m\) with \(Z_m\leq b\) is, path by path, the upper one-sided rule for
 the score \(-z\) with upper boundary \(\log(1/B)=-b\). The score \(-z\)
 has mean \(-Ez>0\) and the same variance \(\sigma_z^2\), and
 \(m=(Ez)^2n/(2\sigma_z^2)\) is unchanged. So the source's continuous
 approximation \(f_c\) of \cref{def:p15-se-transform-approximation},
 applied to \(-z\), assigns to \(m\) the density \(f_c\) with
 \[
  c=\frac{-Ez}{\sigma_z^2}\log\frac1B=\frac{Ez}{\sigma_z^2}\log B>0,
 \]
 which is the source's substitution of \((Ez/\sigma_z^2)\log B\) for
 \(c\) \cite[Section 4.4, p.~150]{Wald1945SequentialTests}. By
 \cref{thm:p15-se-continuous-density} it is a probability density with
 mean \(c/2\) and variance \(c/4\); like the upper case, it is an
 approximation, not the exact law of the integer \(n\).
\item (A lower boundary that is rarely reached.) Let \(T_A\) be the first
 \(m\) with \(Z_m\geq a\), possibly infinite, and \(T\) the two-sided
 exit time. Then \(T=T_A\) on \(Q_1\), so for every set \(G\) of
 values,
 \(|P(T\in G)-P(T_A\in G)|\leq P(T\neq T_A)\leq1-\gamma\). The law of
 the number of observations of the rule with \(B=0\) is therefore
 within \(1-\gamma\) in total variation of the law with \(B>0\). This
 is the exact form of the source's statement that (4.51) remains a good
 approximation for \(B>0\) when \(P(Z_n\geq\log A)\) is nearly one
 \cite[Section 4.4, p.~151]{Wald1945SequentialTests}; it compares the
 two exact laws, and adds nothing to the accuracy of (4.51) for the
 rule with \(B=0\).
\end{enumerate}
\end{hbtheorem}

\begin{hbwarning}{def:p15-w02c-p15-020-non-normal}{Normal approximations to the law of the number of observations}
The source states, citing its reference [4], that when \(|Ez|\) and
\(\sigma_z\) are small enough the law of \(n\) computed as if \(z\) were
normal approximates the exact law for non-normal \(z\). It explains the
use of a density for the discrete \(n\) by the neglected excess, which
is zero only in the limit where \(Ez\) and \(\sigma_z\) tend to zero
\cite[Section 4.4, p.~151]{Wald1945SequentialTests}. The first statement
is an invariance-principle assertion. This handbook proves neither it
nor any error bound for it; what is proved about the law of \(n\) is
\cref{thm:p15-se-transform}, \cref{thm:p15-se-continuous-density} and
\cref{thm:p15-w02c-p15-020-one-sided}.
\end{hbwarning}

\section{The comparison with the current most powerful test}
\label{sec:p15-w02c-p15-020-current}

\begin{hbtheorem}{thm:p15-w02c-p15-020-current-size}{proved=annex:p15-w02c-p15-020-current-size}{Size of the current procedure}
Let \(P_0\neq P_1\) and \(\alpha,\beta\in(0,1)\). For \(N\geq1\) let
\(\beta_N(\alpha)\) be the smallest second-kind error of a test of size
at most \(\alpha\) based on \(x_1,\ldots,x_N\); by
\cref{thm:p15-neyman-pearson} it is attained by a test that rejects
\(H_0\) when \(p_{1N}/p_{0N}>k\) and randomizes on equality, the region
(1.1) of the source. Then:
\begin{enumerate}
\item \(\beta_{N+1}(\alpha)\leq\beta_N(\alpha)\);
\item \(\beta_N(\alpha)\leq\alpha^{-1/2}\rho^N\), where
 \(\rho=\int\sqrt{f_0f_1}\,d\nu<1\); hence \(\beta_N(\alpha)\to0\);
\item the set of \(N\) with \(\beta_N(\alpha)\leq\beta\) is nonempty and
 contains every integer beyond its least element \(N(\alpha,\beta)\).
\end{enumerate}
The current procedure uses \(N(\alpha,\beta)\) observations and the
region (1.1) of size \(\alpha\)
\cite[Introduction, p.~118; Section 1, p.~122]{Wald1945SequentialTests}.
For the normal hypotheses of \cref{thm:p15-w02c-p15-020-small-shift},
\(\beta_N(\alpha)=\Phi(q_{1-\alpha}-2d\sqrt N)\) and
\(N(\alpha,\beta)=\lceil n_d\rceil\), the formula of
\cref{def:p15-se-normal-saving}.
\end{hbtheorem}

\begin{hbtheorem}{thm:p15-w02c-p15-020-current-efficiency}{proved=annex:p15-w02c-p15-020-current-efficiency}{The efficiency of the current test in the small-shift regime}
In the setting of \cref{thm:p15-w02c-p15-020-small-shift}, let \(e_d\)
be the minimax efficiency (\cref{def:p15-se-efficiency}) of the current
most powerful test with \(N(\alpha,\beta)\) observations. Then
\[
 \limsup_{d\to0}e_d\leq\max\{R_0(\alpha,\beta),R_1(\alpha,\beta)\}.
\]
For \(\alpha=\beta=.05\) the right side is about \(.4897\). Among the
twenty-five pairs with \(\alpha,\beta\in\{.01,.02,.03,.04,.05\}\) it is
below \(1/2\) for seventeen; for the other eight, the pairs
\((.01,.04)\), \((.01,.05)\), \((.02,.05)\), \((.03,.05)\) and their
mirror images, it lies between \(.50\) and \(.53\). This is the exact
content, in this regime, of the source's statement that the efficiency
of the current test is rather low and frequently below one half
\cite[Section 2.3, p.~125]{Wald1945SequentialTests}. The bound uses an
equal-strength competitor obtained from a probability ratio test by
randomizing its terminal decision, which leaves its stopping time
unchanged.
\end{hbtheorem}

\begin{hbexample}{def:p15-w02c-p15-020-table-1}{The full savings table}
Recomputed from the limits \(R_1,R_0\) of
\cref{thm:p15-w02c-p15-020-small-shift}, the percentage savings
\(100(1-R_1)\) when \(H_1\) is true, with rows \(\beta\) and columns
\(\alpha\), are
\begin{center}
\begin{tabular}{c|ccccc}
\(\beta\backslash\alpha\)&.01&.02&.03&.04&.05\\\hline
.01&58.40&60.21&61.40&62.33&63.09\\
.02&53.97&55.71&56.86&57.75&58.50\\
.03&51.04&52.71&53.81&54.67&55.39\\
.04&48.82&50.41&51.48&52.30&52.99\\
.05&47.03&48.55&49.57&50.36&51.03
\end{tabular}
\end{center}
Rounded to whole percentages these are the entries of the source's
Table 1, Panel A, except at \(\beta=.02\), \(\alpha=.05\), where the
value \(58.496\) rounds to \(58\) and the source prints \(59\)
\cite[Section 4.3, Table 1, p.~147]{Wald1945SequentialTests}.
Since \(R_0(\alpha,\beta)=R_1(\beta,\alpha)\) identically, the savings
when \(H_0\) is true are the same table with \(\alpha\) and \(\beta\)
interchanged; this is the source's Panel B and its remark that the
symmetry of the normal law gives it from Panel A
\cite[Section 4.3, p.~147]{Wald1945SequentialTests}. The smallest entry
is \(47.03\), at \(\alpha=.01\), \(\beta=.05\), so on the tabulated
pairs the saving is at least \(47\) percent, as the source states for
the range from \(.01\) to \(.05\)
\cite[Section 4.3, p.~148]{Wald1945SequentialTests}. A numerical
evaluation on the grid of step \(.001\) in \([.01,.05]^2\) gives the same
minimum at the same corner; between grid points the bound is not proved
here. The values are limits as \(\theta_1-\theta_0\to0\), with the
targets in place of the actual errors; at a positive shift the exact
ratios differ from the overshoot-free ones by at most
\(2M_d/(q_{1-\alpha}+q_{1-\beta})^2\), by
\cref{thm:p15-w02c-p15-020-small-shift}, item (1).
\end{hbexample}

\section{Truncation}
\label{sec:p15-w02c-p15-020-truncation}

\begin{hbtheorem}{thm:p15-w02c-p15-020-truncation-exact}{proved=annex:p15-w02c-p15-020-truncation-exact}{Exact error change under truncation}
For the finite-score test with \(T<\infty\) almost surely under both
hypotheses, truncate at \(n_0\) as in \cref{thm:p15-se-truncation}:
accept \(H_1\) at \(n_0\) when \(Z_{n_0}>0\) and \(H_0\) otherwise, and
let \(\alpha(n_0),\beta(n_0)\) be the resulting errors. Then
\[
 \alpha(n_0)=\alpha'+\rho_0(n_0)-\rho_0^-(n_0),\qquad
 \beta(n_0)=\beta'+\rho_1(n_0)-\rho_1^-(n_0),
\]
where, with \(Q_i\) the decisions of the untruncated test,
\(\rho_0(n_0)=P_0(T>n_0,Z_{n_0}>0,Q_0)\),
\(\rho_0^-(n_0)=P_0(T>n_0,Z_{n_0}\leq0,Q_1)\),
\(\rho_1(n_0)=P_1(T>n_0,Z_{n_0}\leq0,Q_1)\) and
\(\rho_1^-(n_0)=P_1(T>n_0,Z_{n_0}>0,Q_0)\). The source's bound
\(\alpha(n_0)\leq\alpha+\bar\rho_0(n_0)\), with
\(\bar\rho_0(n_0)=P_0(0<Z_{n_0}<a)\), therefore exceeds the true value
by \(\alpha-\alpha'\) plus
\[
 \rho_0^-(n_0)+P_0(T\leq n_0,0<Z_{n_0}<a)+P_0(T>n_0,Z_{n_0}>0,Q_1),
\]
and the bound \(\beta(n_0)\leq\beta+P_1(b<Z_{n_0}\leq0)\) exceeds
\(\beta(n_0)\) by \(\beta-\beta'\) plus
\(\rho_1^-(n_0)+P_1(T\leq n_0,b<Z_{n_0}\leq0)+P_1(T>n_0,Z_{n_0}\leq0,Q_0)\).
None of these terms need be small.
\end{hbtheorem}

\begin{hbexample}{def:p15-w02c-p15-020-tables-2-3}{The full termination and truncation tables}
Take \(H_0:N(0,1)\) against \(H_1:N(\theta,1)\), with \(\theta\) chosen
so that the current test of strength \((\alpha,\beta)\) needs exactly
\(1000\) observations, and \(a=\log\{(1-\beta)/\alpha\}\),
\(b=\log\{\beta/(1-\alpha)\}\), so \(\mu_1=-\mu_0=\theta^2/2\) and the
score has standard deviation \(\theta\). The table gives, to three
places, the lower bounds \(1-\Phi((a-n_0\mu_1)/(\theta\sqrt{n_0}))\)
under \(H_1\) (alt.) and \(\Phi((b-n_0\mu_0)/(\theta\sqrt{n_0}))\) under
\(H_0\) (null) of \cref{thm:p15-se-deadline}, and the truncation bounds
of \cref{thm:p15-se-truncation} on \(\alpha(n_0)\) and \(\beta(n_0)\)
with the targets as base errors.
Termination lower bounds:
\begin{center}\small
\begin{tabular}{r|cc|cc|cc}
&\multicolumn{2}{c|}{\(\alpha=\beta=.01\)}&\multicolumn{2}{c|}{\(\alpha=.01,\ \beta=.05\)}&\multicolumn{2}{c}{\(\alpha=\beta=.05\)}\\
\(n_0\)&alt.&null&alt.&null&alt.&null\\\hline
1000&.910&.910&.799&.891&.773&.773\\
1200&.950&.950&.870&.932&.838&.838\\
1400&.972&.972&.916&.957&.883&.883\\
1600&.985&.985&.946&.972&.915&.915\\
1800&.991&.991&.965&.982&.938&.938\\
2000&.995&.995&.977&.989&.955&.955\\
2200&.997&.997&.985&.993&.967&.967\\
2400&.998&.998&.990&.995&.976&.976\\
2600&.999&.999&.994&.997&.982&.982\\
2800&1.000&1.000&.996&.998&.987&.987\\
3000&1.000&1.000&.997&.999&.990&.990
\end{tabular}
\end{center}
Truncation bounds:
\begin{center}\small
\begin{tabular}{r|cc|cc|cc}
&\multicolumn{2}{c|}{\(\alpha=\beta=.01\)}&\multicolumn{2}{c|}{\(\alpha=.01,\ \beta=.05\)}&\multicolumn{2}{c}{\(\alpha=\beta=.05\)}\\
\(n_0\)&\(\alpha(n_0)\)&\(\beta(n_0)\)&\(\alpha(n_0)\)&\(\beta(n_0)\)&\(\alpha(n_0)\)&\(\beta(n_0)\)\\\hline
1000&.020&.020&.033&.070&.094&.094\\
1200&.015&.015&.024&.063&.081&.081\\
1400&.013&.013&.019&.058&.072&.072\\
1600&.012&.012&.016&.055&.066&.066\\
1800&.011&.011&.014&.053&.062&.062\\
2000&.010&.010&.012&.052&.058&.058\\
2200&.010&.010&.012&.051&.056&.056\\
2400&.010&.010&.011&.051&.055&.055\\
2600&.010&.010&.011&.051&.053&.053\\
2800&.010&.010&.010&.050&.052&.052\\
3000&.010&.010&.010&.050&.052&.052
\end{tabular}
\end{center}
These agree with the source's Tables 2 and 3 in every entry except
eleven, each differing by one unit in the third place. In Table 2 the source
prints \(.999\) (both columns at \(\alpha=\beta=.01\), \(n_0=2400\)),
\(.871\) (alternative, \(\alpha=.01\), \(\beta=.05\), \(n_0=1200\)) and
\(.837\) (both columns at \(\alpha=\beta=.05\), \(n_0=1200\)) where the
values here are \(.99849\), \(.87040\) and \(.83764\). In Table 3, at
\(\alpha=\beta=.05\) and in both columns, it prints \(.095\), \(.082\)
and \(.053\) at
\(n_0=1000\), \(1200\) and \(2800\), where the values here are
\(.09446\), \(.08137\) and \(.05245\)
\cite[Section 4.5, Table 2, p.~152; Section 4.6, Table 3, p.~155]{Wald1945SequentialTests}.
The entries of Table 3 are bounds for the effective errors, not the
errors themselves. The source states that the bounds (4.69)--(4.70) may
considerably exceed \(\alpha(n_0)\) and \(\beta(n_0)\) and that closer
limits would be desirable; it adds, as Wald's own view, that they
are considerably above the true values when \(n_0\) is not much larger
than the number of observations of the current test
\cite[Section 4.6, p.~154]{Wald1945SequentialTests}.
\Cref{thm:p15-w02c-p15-020-truncation-exact} identifies the terms by
which the bounds exceed the true values; their sizes are not computed
here, so the source's view of their magnitude remains a judgement.
\end{hbexample}

\index{sequential testing!boundary calibration}\index{likelihood ratio!sequential boundaries}
